\section{Variance và std deviation [Trung bình]}
\subsection{Phân tích \& Thiết kế (M0)}
    \begin{itemize}
    \item \textbf{Input:}
        \begin{itemize}
            \item File \texttt{array.txt} chứa:
            \begin{itemize}
                \item Số nguyên $n$ (số lượng phần tử).
                \item $n$ số thực.
            \end{itemize}
        \end{itemize}

    \item \textbf{Xử lý:}
        \begin{itemize}
            \item \textbf{Case biên:}
            \begin{itemize}
                \item $n = 0 \rightarrow$ Ghi ra file chữ \texttt{EMPTY}.
                \item $n = 1 \rightarrow$ $\text{variance} = 0$, $\text{std} = 0$.
                \item Tất cả các phần tử bằng nhau $\rightarrow$ $\text{variance} = 0$, $\text{std} = 0$.
            \end{itemize}
            
            \item \textbf{Module 1 (Đọc mảng):}
            \begin{itemize}
                \item Hàm \texttt{readArray(path, vector<double>\&)}
                \item Bắt lỗi mở file (\text{fail} $\rightarrow$ \texttt{cerr} + \texttt{return false}).
            \end{itemize}

            \item \textbf{Module 2 (Tính trung bình - Mean):}
            \begin{itemize}
                \item Hàm \texttt{calcMean(vector<double>)}
                \item Dùng kiểu \texttt{long double} để cộng dồn tránh sai số.
            \end{itemize}
            
            \item \textbf{Module 3 (Tính phương sai và độ lệch chuẩn):}
            \begin{itemize}
                \item Hàm \texttt{calcVariance(vector<double>, double mean)}
                \item Tính phương sai quần thể (chia cho $n$).
                \item Tích lũy tổng bình phương độ lệch $\sum(x_i - \text{mean})^2$ bằng \texttt{long double} để tránh hiện tượng \textit{catastrophic cancellation} (sai số triệt tiêu).
                \item Độ lệch chuẩn (std deviation) $\sigma = \sqrt{\sigma^2}$. Sử dụng hàm \texttt{sqrt()} trong thư viện \texttt{<cmath>}.
            \end{itemize}

            \item \textbf{Module 4 (Ghi output):}
            \begin{itemize}
                \item Hàm \texttt{writeOutput(path, vector<double>)}
                \item Định dạng số thực bằng \texttt{fixed} và \texttt{setprecision(6)}.
            \end{itemize}
        \end{itemize}

        \item \textbf{Output:}
        File \texttt{output.txt} với định dạng:
        \begin{lstlisting}
            mean=<f> / variance=<f> / std=<f>\end{lstlisting} 
    \end{itemize}

\subsection{Mã nguồn C++11 (Tách module)}
    \begin{lstlisting}
    #include <iostream>
    #include <fstream>
    #include <vector>
    #include <iomanip>
    #include <cmath>

    using namespace std;

    // MODULE 1: Doc mang tu file array.txt
    bool readArray(string path, vector<double>& a) {
        ifstream file(path);
        if (!file.is_open()) {
            cerr << "Khong mo duoc file " << path << "\n";
            return false;
        }

        int n;
        if (file >> n) {
            double x;
            for (int i = 0; i < n; i++) {
                file >> x;
                a.push_back(x);
            }
        }

        file.close();
        return true;
    }

    // MODULE 2: Tinh mean (trung binh)
    double calcMean(const vector<double>& a) {
        if (a.empty()) return 0.0;
        
        long double sum = 0.0;
        for (double x : a) {
            sum += x;
        }
        return (double)(sum / a.size());
    }

    // MODULE 3: Tinh variance (phuong sai quan the)
    double calcVariance(const vector<double>& a, double mean) {
        if (a.size() <= 1) return 0.0;
        
        long double sum_sq_diff = 0.0;
        for (double x : a) {
            double diff = x - mean;
            // Dung long double tich luy de tranh catastrophic cancellation
            sum_sq_diff += (long double)(diff * diff); 
        }
        
        // Phuong sai quan the chia cho n
        return (double)(sum_sq_diff / a.size());
    }

    // MODULE 4: Ghi ket qua ra output.txt
    bool writeOutput(string path, const vector<double>& a) {
        ofstream file(path);
        if (!file.is_open()) {
            cerr << "Khong ghi duoc file " << path << "\n";
            return false;
        }

        if (a.empty()) {
            file << "EMPTY";
            file.close();
            return true;
        }

        double mean = calcMean(a);
        double variance = calcVariance(a, mean);
        double std_dev = sqrt(variance);

        file << fixed << setprecision(6);
        file << "mean=" << mean 
            << " / variance=" << variance 
            << " / std=" << std_dev;

        file.close();
        return true;
    }

    // MAIN
    int main() {
        vector<double> a;

        if (!readArray("array.txt", a)) {
            return 1;
        }

        if (!writeOutput("output.txt", a)) {
            return 1;
        }

        return 0;
    }
    \end{lstlisting}

\subsection{Kế hoạch kiểm thử (Test Cases)}

    \textbf{Test Case 1 (Thông thường)}
    
    \textbf{Input (array.txt):}
    \begin{lstlisting}
    5 10.0 20.0 30.0 40.0 50.0\end{lstlisting}
    \textbf{Expected Output (output.txt):}
    \begin{lstlisting}
    mean=30.000000 / variance=200.000000 / std=14.142136\end{lstlisting}

    \vspace{0.5cm}

    \textbf{Test Case 2 (Case biên -- n=0)}
    
    \textbf{Input (array.txt):}
    \begin{lstlisting}
    0\end{lstlisting}
    \textbf{Expected Output (output.txt):}
    \begin{lstlisting}
    EMPTY\end{lstlisting}

    \vspace{0.5cm}
    
    \textbf{Test Case 3 (Case biên -- n=1)}

    \textbf{Input (array.txt):}
    \begin{lstlisting}
    1 42.5\end{lstlisting}
    \textbf{Expected Output (output.txt):}
    \begin{lstlisting}
    mean=42.500000 / variance=0.000000 / std=0.000000\end{lstlisting}

    \vspace{0.5cm}
    
    \textbf{Test Case 4 (Case biên -- Tất cả phần tử bằng nhau)}

    \textbf{Input (array.txt):}
    \begin{lstlisting}
    4 7.2 7.2 7.2 7.2\end{lstlisting}
    \textbf{Expected Output (output.txt):}
    \begin{lstlisting}
    mean=7.200000 / variance=0.000000 / std=0.000000\end{lstlisting}